% 自 notes/docs/1151/seminar/W09_1109_預習.md（87808f3）由 tools/md_to_tex.py 於 2026-11-09 轉出一次；此後本檔為正本。
\chapter{Kramer〈Judicial Supremacy and the End of Judicial Restraint〉\pg{621 至 634 頁}}
\label{ch:Kramer}

\begin{longtable}[]{@{}
  >{\raggedright\arraybackslash}p{(\linewidth - 4\tabcolsep) * \real{0.1919}}
  >{\raggedright\arraybackslash}p{(\linewidth - 4\tabcolsep) * \real{0.5455}}
  >{\raggedright\arraybackslash}p{(\linewidth - 4\tabcolsep) * \real{0.2626}}@{}}
\caption{Kramer 一文的節次與頁次。}\tabularnewline
\toprule\noalign{}
\begin{minipage}[b]{\linewidth}\raggedright
節次
\end{minipage} & \begin{minipage}[b]{\linewidth}\raggedright
小標
\end{minipage} & \begin{minipage}[b]{\linewidth}\raggedright
頁次
\end{minipage} \\
\midrule\noalign{}
\endfirsthead
\toprule\noalign{}
\begin{minipage}[b]{\linewidth}\raggedright
節次
\end{minipage} & \begin{minipage}[b]{\linewidth}\raggedright
小標
\end{minipage} & \begin{minipage}[b]{\linewidth}\raggedright
頁次
\end{minipage} \\
\midrule\noalign{}
\endhead
\bottomrule\noalign{}
\endlastfoot
甲 & 摘要與論題 & 621--622 \\
乙 & 建國時期的兩種司法審查\ix{terms}{01@\zhlidxcat{司法審查權限之歸屬}!司法審查}理論 & 622--624 \\
丙 & 解釋方法\ix{terms}{05@\zhlidxcat{解釋理論與自制}!解釋理論}與自制\ix{terms}{05@\zhlidxcat{解釋理論與自制}!司法自制} & 624--626 \\
丁 & 1800 年至新政 & 626--630 \\
戊 & Brown 案\ix{cases}{02@美國!01@聯邦最高法院!1954Brown v. Board of Education（1954）@Brown v. Board of Education（1954）}之後的自由派與保守派 & 630--632 \\
己 & 解釋理論的興起與自制的衰退 & 632--634 \\
\end{longtable}

\section{摘要與論題\pg{621 至 622 頁}}
\label{sec:Kramer-1}

\subsection{（一）首頁摘要\pg{621 頁}}

首頁摘要寫道，Thayer\ix{persons}{thayer@Thayer, James Bradley} 重申的是較早的 Jeffersonian\ix{terms}{01@\zhlidxcat{司法審查權限之歸屬}!共和派} 觀念：「primary authority to interpret the Constitution lies with the people and not with courts」〔詮釋憲法的首要權威在人民，不在法院〕。這種 popular constitutionalism\ix{terms}{01@\zhlidxcat{司法審查權限之歸屬}!人民憲政主義}〔人民憲政主義〕被現代的司法至上（judicial supremacy\ix{terms}{01@\zhlidxcat{司法審查權限之歸屬}!司法至上}）取代，後者「took hold only in the 1960s」〔遲至 1960 年代才站穩〕，使政治辯論從誰應詮釋憲法轉為憲法應如何詮釋。司法至上被接受之後，自制\ix{terms}{05@\zhlidxcat{解釋理論與自制}!司法自制}只剩一項理由，即保護司法權威，免於因行動過於頻繁或積極而招致政治反撲；摘要稱這項理由「grows weaker in direct proportion to supremacy\textquotesingle s acceptance」〔隨司法至上被接受的程度而等比減弱〕。

摘要接著引 Posner\ix{persons}{posner@Posner, Richard A.} 的觀察：司法自制今日已退化為「little more than a rhetorical tool used opportunistically by both sides」〔雙方機會主義地使用的修辭工具而已〕。Kramer\ix{persons}{kramer@Kramer, Larry D.} 對成因的看法，摘要中的原文是：「the culprit for this development is not, as Posner suggests, the development of theories of constitutional interpretation」〔此一發展的禍首，並非如 Posner 所言，是憲法解釋理論\ix{terms}{05@\zhlidxcat{解釋理論與自制}!解釋理論}的發展〕，禍首在「our embrace of the notion that deciding the ongoing meaning of a democratic Constitution is a task best left entirely to judges」〔我們接受了這個想法：決定民主憲法持續演變的意義，最宜全交給法官〕。

\subsection{（二）對 Posner 的回應\pg{621 至 622 頁}}

本文開頭以 Posner\ix{persons}{posner@Posner, Richard A.} 與 Holmes\ix{persons}{holmes@Holmes, Oliver Wendell, Jr.（霍姆斯）} 並舉，稱 Posner 之於二十世紀末與二十一世紀初，如 Holmes 之於十九世紀末與二十世紀初\pg{621 頁}。622 頁接著說，Posner 的論述有趣，不等於完全正確；Kramer\ix{persons}{kramer@Kramer, Larry D.} 認為 Posner「misunderstood important historical nuances」〔誤解了重要的歷史細微處〕，因此要講「a somewhat different story」〔一個略有不同的故事〕，提出可以並列考慮的另一組因素。被回應的文章是 Posner, \en{The Rise and Fall of Judicial Self-Restraint}, 100 Calif. L. Rev. 519 (2012)，註 8\pg{624 頁}列有出處。首頁註 * 僅載明 Kramer 為 Stanford Law School 的 Richard E. Lang Professor and Dean。

\subsection{（三）自制的興衰與起點\pg{622 頁}}

本文認為自制\ix{terms}{05@\zhlidxcat{解釋理論與自制}!司法自制}「rose but has fallen」〔興起而後衰落〕一事「seems plain enough」〔相當明顯〕。衰落的原因不能歸於左派或右派，只能說無論哪一方掌權，都似乎決意超過對方上一次的作為，因此這是「an accelerating slide of many years」〔多年來不斷加速的下滑〕。至於興起，Posner\ix{persons}{posner@Posner, Richard A.} 追溯到 Thayer\ix{persons}{thayer@Thayer, James Bradley} 1893 年的 \en{The Origin and Scope of the American Doctrine of Constitutional Law}（註 1：7 Harv. L. Rev. 129）；Kramer\ix{persons}{kramer@Kramer, Larry D.} 回答這「much too late」〔太晚了〕，要理解自制作為對政治決策者的尊重（deference\ix{terms}{05@\zhlidxcat{解釋理論與自制}!尊重}）之興衰，須回到建國時期與司法審查\ix{terms}{01@\zhlidxcat{司法審查權限之歸屬}!司法審查}的起源，Thayer 自己也是如此。司法審查如何發展，Kramer 指向自己的 \en{The People Themselves}（2004）第 9 至 144 頁，另列 Friedman\ix{persons}{friedman@Friedman, Barry} 與 Hamburger 兩部近著（註 2）。

\section{建國時期的兩種司法審查理論\pg{622 至 624 頁}}
\label{sec:Kramer-2}

\subsection{（一）共和派：人民憲政主義\pg{622 至 623 頁，註 3 至 5}}

Kramer\ix{persons}{kramer@Kramer, Larry D.} 稱建國時期出現兩套正當化司法審查\ix{terms}{01@\zhlidxcat{司法審查權限之歸屬}!司法審查}的理由，並稱此二說之出現無爭議。第一套稱為「Republican」或「Jeffersonian\ix{terms}{01@\zhlidxcat{司法審查權限之歸屬}!共和派}」，因 Jefferson\ix{persons}{jefferson@Jefferson, Thomas（傑佛遜）} 一派所倡而得名，出現於 1780 年代，前提是「primary authority to interpret, as well as to make, constitutional law rested actively in the community」〔詮釋與制定憲法的首要權威都實際落在社群手中〕，即他在別處所稱的 popular constitutionalism\ix{terms}{01@\zhlidxcat{司法審查權限之歸屬}!人民憲政主義}\pg{註 3，引 \en{The People Themselves} 第 8 頁}\pg{622 頁}。

依這種看法，憲法是人民制定來規範並約束政府的法，普通法則是政府制定來規範並約束人民的法，這個倒置又倒置了解釋與執行權限的通常分配\pg{622 頁}。政府官員是普通法的權威詮釋者；人民只能透過選舉與罷免對立法有間接控制，法律一旦制定，公民在確定其意義與施加制裁上居於官員之下，自己的詮釋只是「guesses」〔臆測〕，沒有正式法律權威\pg{622 至 623 頁}。

憲法的關係則相反。Kramer 引共和派的 William Nelson 法官在 \en{Kamper v. Hawkins\ix{cases}{02@美國!02@維吉尼亞州一般法院!1793Kamper v. Hawkins（1793）@Kamper v. Hawkins（1793）}}, 3 Va. (1 Va. Cas.) 20, 24 (1793) 的話：「{[}a{]} Constitution . . . is to the governors, or rather to the departments of government, what a law is to individuals」〔憲法之於統治者，或毋寧說之於政府各部門，如同法律之於個人〕\pg{623 頁，註 4}。憲法的對象是公職人員，官員須盡力詮釋並遵循憲法，他們的詮釋與人民對普通法的詮釋一樣不具權威，受「the People themselves」〔人民自己〕這一更高權威的直接監督與矯正\pg{623 頁，註 5}。

\subsection{（二）聯邦派：司法至上的先聲\pg{623 頁，註 6}}

第二套理論出現在 1790 年代，憲法批准之後，由 Federalist 領袖發展。Kramer\ix{persons}{kramer@Kramer, Larry D.} 把它放在更大的思想演變中：憲法開始運作，法國大革命爆發，Federalist 為了在看來民主失控的局面中確保「order」〔秩序〕，重新思考司法權的性質，形成「what amounts to the modern theory of judicial supremacy\ix{terms}{01@\zhlidxcat{司法審查權限之歸屬}!司法至上}」〔相當於現代司法至上理論的主張〕\pg{623 頁}。

這套理論把憲法視為「nothing more than a species of ordinary law」〔不過是普通法的一種〕，憲法解釋因此是普通的司法與法律行為，不授權給立法或行政官員，更不授權給全體人民。政治部門行動前仍須自行詮釋憲法，依權力分立\ix{terms}{06@\zhlidxcat{審查制度之類型}!權力分立體制}，首要的解釋權在法院，法院對憲法意義與合憲性有「special and final responsibility」〔特別且終局的責任〕。Kramer 以 James Kent 1794 年的話為例：憲法的解釋或建構「requires the exercise of the same LEGAL DISCRETION as the interpretation or construction of a Law」〔需要與解釋或建構一部法律相同的法律裁量〕，所以法官是「the proper and intended Guardians of our limited Constitutions」〔有限憲法適當而原定的守護者〕\pg{623 頁，註 6}。

\subsection{（三）羅德島議會的答覆\pg{624 頁，註 7}}

624 頁開頭，羅德島議會拒絕 Madison\ix{persons}{madison@Madison, James（麥迪遜）} 集結人民反對 Alien and Sedition Acts 之合憲性的努力，說得「still more baldly」〔更為直白〕：憲法「vests in the federal courts, exclusively, and in the Supreme Court of the United States, ultimately, the authority of deciding on the constitutionality of any act or law of the Congress of the United States」〔把決定國會任何法案或法律合憲性的權限，專屬授予聯邦法院，並以聯邦最高法院為終局〕\pg{註 7，載 Elliot 編 \en{Debates} 第 4 卷 532, 533 頁}。

\section{解釋方法與自制\pg{624 至 626 頁}}
\label{sec:Kramer-3}

\subsection{（一）雙方都未爭論如何解釋憲法\pg{624 頁}}

Kramer\ix{persons}{kramer@Kramer, Larry D.} 就兩派的差異提出兩點。第一點，雙方都沒有花時間討論如何解釋憲法，也沒有興趣發展今日意義的憲法解釋「theory」〔理論〕，因為對這個問題並無真正的分歧。大家基本上相信，憲法可以而且應當以橫跨各法領域的同一種開放的論辯程序來解釋，即相當不拘結構地運用文本、結構、歷史、先例與後果的論據，得出最有說服力的整體結論\pg{624 頁}。

既然對如何解釋看法一致，爭議就集中於誰能解釋、發生什麼效果。Kramer 補充，解釋過程鬆散而開放，不同的解釋者可能對同一問題合理地得出不同結論，所以誰有解釋權限\ix{terms}{01@\zhlidxcat{司法審查權限之歸屬}!解釋權限}才如此關鍵\pg{624 頁}。

\subsection{（二）自制見於兩套理論，Thayer 並非原創\pg{624 頁，註 8}}

第二點與本文主題更直接相關：自制\ix{terms}{05@\zhlidxcat{解釋理論與自制}!司法自制}的概念在兩套司法審查\ix{terms}{01@\zhlidxcat{司法審查權限之歸屬}!司法審查}理論裡都有位置，並且「central to both approaches」〔對兩者都居核心〕。Kramer\ix{persons}{kramer@Kramer, Larry D.} 因此不同意 Posner\ix{persons}{posner@Posner, Richard A.} 的說法：Thayer\ix{persons}{thayer@Thayer, James Bradley} 所憑藉的不只是「scattered judicial remarks」〔零散的司法評論〕，在「organizing and rationalizing」〔整理並使之合理化〕它們上也不是首創\pg{註 8，引 Posner 523 頁}。依 Kramer，Thayer 是要把一個「old and well-developed concept of self-restraint」〔古老而發展成熟的自制概念〕重新推回舞台中央，當時法院及其支持者正傾向忽視它，情形與今日相仿\pg{624 頁}。

\subsection{（三）Posner 的「無理論」前提\pg{624 至 625 頁，註 9}}

Kramer\ix{persons}{kramer@Kramer, Larry D.} 不同意 Posner\ix{persons}{posner@Posner, Richard A.} 認為自制\ix{terms}{05@\zhlidxcat{解釋理論與自制}!司法自制}的前提是「hav{[}ing{]} no theory」〔沒有理論〕、自制之死可歸因於宣稱能產生「correct」〔正確〕憲法決定的解釋理論\ix{terms}{05@\zhlidxcat{解釋理論與自制}!解釋理論}之興起\pg{624 頁，註 9，引 Posner 538 頁}。理由是即使主張只有單一正解的理論，適用上也有足夠的不確定性，使理性的人對何者為正解合理地有歧見，Kramer 說這就足以為尊重\ix{terms}{05@\zhlidxcat{解釋理論與自制}!尊重}或自制創造空間。他並指出十八、十九世紀的律師也相信自己的方法有單一正解\pg{625 頁}。

Kramer 同意 Posner 的是，今日法院與評論者都不太傾向自制。他的主張是這項變化另有來源，解釋理論的發展本身，是削弱自制的同一組發展的「by-product」〔副產品〕，後文再回到這個問題\pg{625 頁}。

\subsection{（四）共和派的自制理由\pg{625 至 626 頁，註 10、11}}

兩派都期待法院在行使司法審查\ix{terms}{01@\zhlidxcat{司法審查權限之歸屬}!司法審查}時自制\ix{terms}{05@\zhlidxcat{解釋理論與自制}!司法自制}，理由與目的不同\pg{625 頁}。對共和派\ix{terms}{01@\zhlidxcat{司法審查權限之歸屬}!共和派}，自制是保存「the People」首要且優越之權威的手段。人民既已把制定法律的責任委託給立法機關，一項法律是否合憲，首先是社群與民選代表之間透過抗議、請願、選舉與其他直接的人民行動來解決的事。少數共和派因此主張任何形式的司法審查都不正當\pg{註 10，引 11 Annals of Cong. 178--80 (1802) Breckinridge 參議員之發言，及同冊 661 頁 Randolph 眾議員之發言}。

共和派的正統立場承認法院有時必須拒絕執行被認為違憲的法律。人民既然委託法院裁判案件，法院在裁判時若不考慮合憲性，就違反了作為人民代理人的責任，與立法者不審酌法律是否合憲同理。法院宣告立法無效時，是以人民代理人的身分行動。遇有疑義，只有該法「unconstitutional beyond dispute」〔違憲而無可爭辯〕時才能宣告無效\pg{註 11，引 Iredell\ix{persons}{iredell@Iredell, James} 1787 年 8 月 26 日致 Spaight 函；並載 Thayer\ix{persons}{thayer@Thayer, James Bradley} 138--42 頁彙整之早期案例，及 Treanor 58 Stan. L. Rev. 455 (2005)}。Kramer\ix{persons}{kramer@Kramer, Larry D.} 說，代理人揣度本人（principal\ix{terms}{01@\zhlidxcat{司法審查權限之歸屬}!本人與代理人}）的指令，不可輕易為之，須待近於確定的條件，因為本人自己能夠行動，且始終保有首要責任\pg{625 至 626 頁}。

\subsection{（五）聯邦派的自制理由\pg{626 頁，註 12 至 16}}

聯邦派\ix{terms}{01@\zhlidxcat{司法審查權限之歸屬}!聯邦派}明確否認人民在詮釋憲法上有首要的、乃至任何的權威（相對於制定新的憲法法）。依 Kramer\ix{persons}{kramer@Kramer, Larry D.}，到 1790 年代中期，聯邦派採取的立場已與今日多數人相近：憲法把終局的解釋權限\ix{terms}{01@\zhlidxcat{司法審查權限之歸屬}!解釋權限}授予法院。法官解釋憲法時不是社群的代理人\ix{terms}{01@\zhlidxcat{司法審查權限之歸屬}!本人與代理人}；他們在阻止社群及其立法與行政代理人犯下 Gouverneur Morris\ix{persons}{morris@Morris, Gouverneur} 所說的「the greatest follies and absurdities」〔最大的愚行與荒謬〕（註 12）。Morris 稱讓人民解釋憲法，就是把「an inestimable treasure . . . into the hands of drunkards, madmen, and fools」〔無價之寶\ldots\ldots 交到醉漢、瘋子與傻瓜手中〕（註 13）。因此憲法刻意「vested in the judges a check」〔把一項制衡授予法官〕（註 14），照 Morris 的嘲諷，是為了「to save the people from their most dangerous enemy, to save them from themselves」〔使人民免於其最危險的敵人，即免於自己〕（註 15）。

聯邦派同樣強調，法律的違憲須屬明確才能不予執行，「in all doubtful cases . . . the Act ought to be supported」〔凡有疑義之案\ldots\ldots 該法應予維持〕\pg{註 16，引 Iredell\ix{persons}{iredell@Iredell, James} 函第 175 頁}。但在聯邦派，這種自制\ix{terms}{05@\zhlidxcat{解釋理論與自制}!司法自制}是審慎\ix{terms}{05@\zhlidxcat{解釋理論與自制}!審慎}與政治權宜的問題，目的是藉由降低越權的風險，保障並維持司法權威，原文括號標明此為 judicial (rather than popular) authority\pg{626 頁}。

\section{1800 年至新政\pg{626 至 630 頁}}
\label{sec:Kramer-4}

\subsection{（一）1800 年選舉與 Marbury\pg{626 頁，註 17、18}}

Kramer\ix{persons}{kramer@Kramer, Larry D.} 寫道，1800 年的選舉決定性地解決了兩種觀點的爭議，由共和派\ix{terms}{01@\zhlidxcat{司法審查權限之歸屬}!共和派}勝出\pg{註 17，引 \en{The People Themselves} 137--38 頁}。\footcite{kramer2004}聯邦派\ix{terms}{01@\zhlidxcat{司法審查權限之歸屬}!聯邦派}又延續幾年，在 1812 年戰爭之後完全消失；John Marshall\ix{persons}{marshall@Marshall, John（馬歇爾）} 在 \en{Marbury\ix{cases}{02@美國!01@聯邦最高法院!1803Marbury v. Madison（1803）@Marbury v. Madison（1803）} v. Madison\ix{persons}{madison@Madison, James（麥迪遜）}} 中明顯放棄聯邦派的司法審查\ix{terms}{01@\zhlidxcat{司法審查權限之歸屬}!司法審查}觀，所作意見是 Jefferson\ix{persons}{jefferson@Jefferson, Thomas（傑佛遜）} 自己可能寫出的，且大部分取自維吉尼亞共和派 Spencer Roane 與 St. George Tucker 在 \en{Kamper v. Hawkins\ix{cases}{02@美國!02@維吉尼亞州一般法院!1793Kamper v. Hawkins（1793）@Kamper v. Hawkins（1793）}} 的意見\pg{626 頁，註 18，另引 Kramer, \en{Marbury and the Retreat from Judicial Supremacy}, 20 Const. Comment. 205, 228--29 (2003)}。

\subsection{（二）聯邦派觀念的殘存\pg{626 至 627 頁，註 19、20}}

選舉很少殺死觀念，所以聯邦派\ix{terms}{01@\zhlidxcat{司法審查權限之歸屬}!聯邦派}的司法審查\ix{terms}{01@\zhlidxcat{司法審查權限之歸屬}!司法審查}概念並未完全消失，存續於一些前聯邦派人士（後來成為 National Republicans，再成為 Whigs）之間，尤其是逐漸職業化的法官與律師\pg{626 頁，註 19，引 \en{The People Themselves} 156--64 頁}。法院對憲法有特別權威、優越於人民與其他部門的想法，因此一直留在國家的政治與法律論述中；每當其倡議者推得夠用力、引起公眾注意，就被再度壓下，Kramer\ix{persons}{kramer@Kramer, Larry D.} 舉 Second Bank of the United States 之爭與 \en{Dred Scott\ix{cases}{02@美國!01@聯邦最高法院!1857Dred Scott v. Sandford（1857）@Dred Scott v. Sandford（1857）}} 引發的爭議為例\pg{627 頁，註 20}。

\subsection{（三）重建之後司法至上再度出現\pg{627 頁，註 21 至 25}}

Kramer\ix{persons}{kramer@Kramer, Larry D.} 列出重建之後司法至上\ix{terms}{01@\zhlidxcat{司法審查權限之歸屬}!司法至上}強勢重現的幾項因素：\en{Dred Scott\ix{cases}{02@美國!01@聯邦最高法院!1857Dred Scott v. Sandford（1857）@Dred Scott v. Sandford（1857）}} 的重量隨時間與法院人事更替而逐漸減輕；共和黨的優勢在 1874 年後讓位給數十年的分裂政府，使大法官較自由行動（註 21，引 Lichtman \& DeCell；Gely \& Spiller, 12 Int\textquotesingle l Rev. L. \& Econ. 45 (1992)）；內戰時期的自由勞動修辭，在保守派手中演變為新的自由放任、反重分配的正統（註 22，引 Foner；Forbath, 1985 Wis. L. Rev. 767）。法院在道德上忠於特定的財產權觀念，又懼怕勞工運動與其他社會動盪的徵兆，開始擴張自己的角色。Brewer 大法官 1893 年的演講提到「{[}m{]}agnifying, like the apostle of old, my office」〔像古代使徒一樣，頌揚我的職務〕，理由是「the salvation of the nation . . . rests upon the independence and vigor of the judiciary」〔國家的救贖\ldots\ldots 繫於司法的獨立與活力〕（註 23）。此時也再度聽到被棄的聯邦派\ix{terms}{01@\zhlidxcat{司法審查權限之歸屬}!聯邦派}舊論，如「what is now to be feared and guarded against is the despotism of the many---of the majority」〔現在該恐懼與防範的，是眾人的專制，即多數的專制〕，以及藉司法審查\ix{terms}{01@\zhlidxcat{司法審查權限之歸屬}!司法審查}「people had effectually protected themselves against themselves」〔人民已有效地保護自己免於自己〕（註 24，引 Dillon 1892 年 ABA 會長演說）。

這些發展逐漸改變法院的憲法判決。因變化緩慢，且起初主要影響已失去強烈政治支持的法律，法院過了幾年才開始受到抨擊；到 1880 年代末，法律評論者與政治領袖已開始注意，一場長期的鬥爭隨之展開，最後在 1930 年代的新政危機（New Deal crisis）告終（註 25，引 Friedman\ix{persons}{friedman@Friedman, Barry}, 76 N.Y.U. L. Rev. 1383 (2001)）。

\subsection{（四）進步時代的爭議\pg{627 至 628 頁，註 26}}

Kramer\ix{persons}{kramer@Kramer, Larry D.} 說，今日的法律人、法官與學者常以為法院的至上地位在這段期間沒有受到挑戰，這是照單全收司法意見修辭所造成的「historical blind spot」〔歷史盲點〕\pg{627 至 628 頁}。大法官關於司法在憲法體制中位置的陳述，「must be seen for what they were: partisan claims in contested territory」〔須視為其本來面目：在爭議地帶的黨派主張〕。進步時代是「flux and uncertainty」〔變動與不確定〕的時期，法院的角色並非核心議題，但法院因積極地在初現的階級衝突中選邊，使自己成為爭議來源，它為自己創造的角色是否適當，成為懸而未決的問題之一\pg{628 頁}。

\subsection{（五）Thayer 的主張\pg{628 頁，註 26 至 28}}

Thayer\ix{persons}{thayer@Thayer, James Bradley} 的文章寫於這場辯論之中，他是眾多提出類似論點者之一（註 26，引 Friedman\ix{persons}{friedman@Friedman, Barry}；T.W. Brown；W.F. Dodd；Meigs 等）。Kramer\ix{persons}{kramer@Kramer, Larry D.} 認為他的特定論點是重申並恢復 Jeffersonian\ix{terms}{01@\zhlidxcat{司法審查權限之歸屬}!共和派} 司法權威觀的優先地位。

Posner\ix{persons}{posner@Posner, Richard A.} 說 Thayer 最大的憂慮是：法院若把憲法限制執行到底，立法者將不再思考擬議立法的合憲性，只想法院會如何反應\pg{註 27，引 Posner 524 頁}。Kramer 說 Thayer 確實提過這個論點，但僅在文章最後一段帶過；把它視為 Thayer 的主要關切，是「read him anachronistically」〔以時代錯置的方式閱讀他〕。Kramer 的讀法是：從文章結構與 Thayer 重複該點的次數可見，他的主要關切是重申詮釋憲法的首要權威在法院之外，宣告法律違憲的司法權威至多是「subordinate, secondary check」〔從屬的、次要的制衡〕\pg{註 28，引 Thayer 第 136、144、148、150、152 頁}。Thayer 要恢復的是共和派較古老、歷史上居首的司法權威觀，包括其自制\ix{terms}{05@\zhlidxcat{解釋理論與自制}!司法自制}與尊重\ix{terms}{05@\zhlidxcat{解釋理論與自制}!尊重}的主張，並拒絕鍍金時代法院對憲法至上的僭越。

\subsection{（六）Thayer 對 Jeffersonian 理論的改動\pg{628 至 629 頁}}

Kramer\ix{persons}{kramer@Kramer, Larry D.} 承認 Thayer\ix{persons}{thayer@Thayer, James Bradley} 手中的 Jeffersonian\ix{terms}{01@\zhlidxcat{司法審查權限之歸屬}!共和派} 理論有所改變。最主要的是，早期共和國所假定且廣被接受的人民直接監督，已不再說得通；Thayer 改以政治部門，尤其立法機關，為人民意見的最佳表達\pg{628 頁}。「The crucial point」〔關鍵〕與 Thayer 的主要對象，是否認憲法的解釋權首先授予法院。Thayer 主張的真相正好相反，立法機關對合憲性的判斷應受尊重\ix{terms}{05@\zhlidxcat{解釋理論與自制}!尊重}，法院的角色至多是在立法機關明顯逾越理性之人所能說是允許的界限時介入，此為最低限度的角色\pg{628 至 629 頁}。

\subsection{（七）對立法機關能力的保留\pg{629 頁，註 29、30}}

Kramer\ix{persons}{kramer@Kramer, Larry D.} 另就 Posner\ix{persons}{posner@Posner, Richard A.} 把 Thayer\ix{persons}{thayer@Thayer, James Bradley} 對立法機關的信賴斥為天真，提出兩點保留，並說這個題目足以另文討論\pg{629 頁}。

\begin{enumerate}
\def\labelenumi{\arabic{enumi}.}
\tightlist
\item
  立法機關無力認真對待憲法這項假定，如同 Posner 的處理所示，比司法至上\ix{terms}{01@\zhlidxcat{司法審查權限之歸屬}!司法至上}更被視為理所當然，也更少被法律學者檢驗\pg{註 29，引 \en{The People Themselves} 238--39 頁；Whittington, 80 N.C. L. Rev. 773, 819--22 (2002)；Burgess；Morgan}。Kramer 稱幕僚與委員會花數百小時理解、辯論並解決憲法問題，所用時間比法院多、資訊比法院好；法院只花一小時聽言詞辯論，評議時間更短，實際審議甚少，多數工作由一名沒有經驗的助理在僅有黨派書狀的情形下完成。高層次的議場辯論或許已成過去，但政治學文獻顯示，認真的審議在委員會、黨團、個人之間、電子郵件與幕僚中進行（註 30，引 Tushnet\ix{persons}{tushnet@Tushnet, Mark}, \en{Taking the Constitution Away from the Courts} 65--66 (1999) 等）。
\item
  立法者會為達政治上想要的結果而行事，不能因此否定這些工作的成果；政治固然居首要地位，國會審議涉及的遠不止於政治。Kramer 也不認為典型最高法院大法官比典型國會議員更不會操縱結果，不明白何以迎合意識形態（他說這顯然驅動大法官的決定）比迎合政治更不令人憂慮，「These days, I\textquotesingle m not even sure they are different」〔如今我甚至不確定兩者有何不同〕\pg{629 頁}。
\end{enumerate}

\subsection{（八）進步派與新政\pg{629 至 630 頁，註 31 至 34}}

回到原來的司法審查\ix{terms}{01@\zhlidxcat{司法審查權限之歸屬}!司法審查}概念：它關切以自制\ix{terms}{05@\zhlidxcat{解釋理論與自制}!司法自制}保存人民對憲法的持續控制，這個概念發生了什麼？Kramer\ix{persons}{kramer@Kramer, Larry D.} 說支持者與司法至上\ix{terms}{01@\zhlidxcat{司法審查權限之歸屬}!司法至上}的支持者的鬥爭延續到二十世紀初幾十年，是 1790 年代共和派\ix{terms}{01@\zhlidxcat{司法審查權限之歸屬}!共和派}與聯邦派\ix{terms}{01@\zhlidxcat{司法審查權限之歸屬}!聯邦派}衝突的一次「close, if prolonged, replay」〔接近而冗長的重演〕\pg{629 至 630 頁}。Thayer\ix{persons}{thayer@Thayer, James Bradley} 與其他進步派（包括常被忽略的 Theodore Roosevelt\ix{persons}{roosevelt@Roosevelt, Franklin D.（羅斯福）}，註 31）大體成功地對抗法院的保守派多數，該多數有時言語大膽，偶爾依其修辭行動，但多半默許州與聯邦各層級管制的穩定成長。1920 年代法院稍為大膽，Kramer 稱 \en{Lochner\ix{cases}{02@美國!01@聯邦最高法院!1905Lochner v. New York（1905）@Lochner v. New York（1905）}} 時代最高點仍最宜描述為法院對新型國家干預與管制的「grudging acceptance」〔勉強接受〕\pg{630 頁，註 32}。

法院因 New Deal 恐慌而試圖劃線時，為時已晚，遭到對抗、威脅並被迫退讓。Kramer 引 Roosevelt 的話：「The Constitution of the United States」是「a layman\textquotesingle s document, not a lawyer\textquotesingle s contract」〔外行人的文件，不是律師的契約〕，其意義與適用應由「the people and the Congress」〔人民與國會〕決定並確定，律師與法官則不在其列\pg{630 頁，註 33，引 1937 年 9 月 17 日憲法日演說}。Roosevelt 隨後藉任命權鞏固勝利，使法院回到願意在當時重要的所有問題上尊重\ix{terms}{05@\zhlidxcat{解釋理論與自制}!尊重}政治部門的狀態\pg{630 頁}。

\section{Brown 案之後的自由派與保守派\pg{630 至 632 頁}}
\label{sec:Kramer-5}

\subsection{（一）Brown 案與司法至上的普遍接受\pg{630 至 631 頁，註 34、35}}

Kramer\ix{persons}{kramer@Kramer, Larry D.} 說轉折「no surprise here」〔不意外〕始於 \en{Brown v. Board of Education\ix{cases}{02@美國!01@聯邦最高法院!1954Brown v. Board of Education（1954）@Brown v. Board of Education（1954）}}。\en{Brown} 之後，司法至上\ix{terms}{01@\zhlidxcat{司法審查權限之歸屬}!司法至上}的觀念似乎終於逐漸獲得廣泛的公眾接受；法院的判決仍常有爭議，但自 1960 年代起，抵抗與不服從轉為抗議的形式，不再主張解釋上的平等，更不用說優越。到 1980 年代，涉及憲法的抗議多半朝向法院（directed to），不再針對法院（against），司法至上的接受成為常態\pg{630 頁，註 34}。法官任命所受注意的急遽增加（註 35，引 Steigerwalt、Wittes、Stras），Kramer 稱為在司法決定之下、任命成為重塑憲法唯一實際途徑的世界中，「a thoroughly predictable development」〔完全可預期的發展〕\pg{630 至 631 頁}。解釋這個「rather extraordinary development」〔相當特殊的發展〕不容易，也在本文範圍之外；Kramer 指出一項明顯重要的來源是 Warren Court 本身，它是美國史上第一次給進步派理由把司法視為朋友\pg{631 頁}。

\subsection{（二）自由派\pg{631 頁，註 36、37}}

自由派知識分子在實質上強烈支持法院所為，但他們的師長與英雄曾領導對抗 \en{Lochner\ix{cases}{02@美國!01@聯邦最高法院!1905Lochner v. New York（1905）@Lochner v. New York（1905）}}，許多人把自己的職業生涯投入這個想法：法院干預人民意志是不當的。\en{Brown} 與隨後其他司法創新，因此是對自由派傳統上承諾司法自制\ix{terms}{05@\zhlidxcat{解釋理論與自制}!司法自制}的「wrenching test」〔痛苦的考驗〕，起初自由派難以決定如何回應 Warren Court（註 36，引 Learned Hand, \en{The Bill of Rights}；Friedman\ix{persons}{friedman@Friedman, Barry}；Peller；Wechsler\ix{persons}{wechsler@Wechsler, Herbert}, \en{Toward Neutral Principles of Constitutional Law}, 73 Harv. L. Rev. 1 (1959)；註 37 另引 Horwitz、Kalman 等）。

Warren Court 的積極作為在 1960 年代中期達到高峰時，新一代自由派學者放棄反對法院，「turned the liberal tradition on its head」〔把自由派傳統倒轉過來〕，接受廣泛司法權威的哲學。這些年輕學者雖熟知 \en{Lochner} 的歷史，卻沒有親歷；對他們而言，當時所稱的 ``the countermajoritarian\ix{terms}{05@\zhlidxcat{解釋理論與自制}!反多數決困難} difficulty''〔反多數決困難〕看來「just did not seem so difficult」〔並不那麼困難〕。他們擱置對課責的疑慮，承認司法審查\ix{terms}{01@\zhlidxcat{司法審查權限之歸屬}!司法審查}可能與民主有緊張關係，並以法院能促進更重要的社會正義為由，證成這項取捨\pg{631 頁}。

\subsection{（三）保守派\pg{631 至 632 頁，註 38 至 40}}

保守派對司法至上\ix{terms}{01@\zhlidxcat{司法審查權限之歸屬}!司法至上}沒有問題。Kramer\ix{persons}{kramer@Kramer, Larry D.} 說這是他們自聯邦派\ix{terms}{01@\zhlidxcat{司法審查權限之歸屬}!聯邦派}時代以來一直接受的想法；他們對 Warren Court 的不滿僅在於其決定是錯的，並抗議自由派大法官以憲法為掩護處理憲法其實不涉及的事項\pg{631 頁}。從 Robert Bork 1968 年在 \en{Fortune} 雜誌對法院的攻擊起（註 38），許多保守派開始使用傳統上屬於自由派的民主與自制\ix{terms}{05@\zhlidxcat{解釋理論與自制}!司法自制}修辭\pg{632 頁}。Kramer 說，他們並未主張最高法院應就何者合憲而尊重\ix{terms}{05@\zhlidxcat{解釋理論與自制}!尊重}政治機關，因為他們支持的是司法至上，所對的是他們所理解與解釋的憲法（註 39，引 Epstein、Macedo、Powell）。他們也一直願意宣告違反原始理解的法律無效，如財產權與槍枝（註 40，引 \en{Kelo v. City of New London\ix{cases}{02@美國!01@聯邦最高法院!2005Kelo v. City of New London（2005）@Kelo v. City of New London（2005）}}, 545 U.S. 469, 505--23 (2005) (Thomas, J., dissenting)；\en{District of Columbia v. Heller\ix{cases}{02@美國!01@聯邦最高法院!2008District of Columbia v. Heller（2008）@District of Columbia v. Heller（2008）}}, 554 U.S. 570 (2008)）。

他們對憲法範圍的看法不同，認為範圍較窄，反對自由派大法官以憲法在其認為憲法未規範的領域制定政策，這些領域因此應屬一般政治。Kramer 結論是，美國史上第一次，保守派與自由派在司法至上的原則上達成一致：法官應當擔任解釋者，且其解釋應為終局且有拘束力\pg{632 頁}。

\section{解釋理論的興起與自制的衰退\pg{632 至 634 頁}}
\label{sec:Kramer-6}

\subsection{（一）爭議的焦點由誰解釋轉為如何解釋\pg{632 至 633 頁}}

此時雙方歧見改為如何解釋。Kramer\ix{persons}{kramer@Kramer, Larry D.} 回到開頭所述：在此之前，並無特別的憲法解釋理論\ix{terms}{05@\zhlidxcat{解釋理論與自制}!解釋理論}，存在的是未經檢驗的共識，即憲法可以也應當像其他法源一樣，用律師從其他脈絡熟悉的、經過檢驗的法律論證方法解釋，至多有些小變化。因為這套制度相對開放，不同解釋者可以用相同的一般方法輕易得出不同結論，所以誰有解釋權限\ix{terms}{01@\zhlidxcat{司法審查權限之歸屬}!解釋權限}很重要；爭議一向在於法官是否或應否從屬於其他解釋者，即立法者、行政首長與社群全體\pg{632 頁}。

既然誰有解釋權已定於法官，如何解釋的問題取得新的重要性，Kramer 稱這「giving rise for the first time to an active focus on interpretive methodology」〔首次產生對解釋方法的積極關注〕，並說明了關於憲法應如何建構之新理論的突然出現。左右兩派的法官與學者忽然發現，憲法須依獨特而特別的規則解釋，而各方產生的規則，似乎必然直接導向該方所偏好的結果\pg{632 至 633 頁}。

\subsection{（二）自制在兩派理論中的位置\pg{633 頁，註 41、42}}

自制\ix{terms}{05@\zhlidxcat{解釋理論與自制}!司法自制}與尊重\ix{terms}{05@\zhlidxcat{解釋理論與自制}!尊重}起初仍是自由派與保守派憲法理論的一環。對自由派，\en{Lochner\ix{cases}{02@美國!01@聯邦最高法院!1905Lochner v. New York（1905）@Lochner v. New York（1905）}} 的教訓延續為有限度審查，但這項教訓現在在實質上限於經濟管制及其他在左派關切的民權、少數者權利與平等保護之外的法律；在後者領域，積極的司法審查\ix{terms}{01@\zhlidxcat{司法審查權限之歸屬}!司法審查}成為「constitutional duty」〔憲法上的義務〕（註 41，引 Dworkin\ix{persons}{dworkin@Dworkin, Ronald（德沃金）}, \en{Taking Rights Seriously} 131--50；Ely, \en{Democracy and Distrust}；Ball, 59 Iowa L. Rev. 1059 (1974)）。對保守派情形正好相反：如前所述，自制成為限制 Warren 與 Burger 早期法院之自由派積極作為的論據，同時他們批評這些法院，使「true」Constitution〔真正的憲法〕處於「exile」〔流放〕，放棄監督國會權力並保護保守派所重視之其他權利的責任（註 42，引 Ginsburg, \en{Delegation Running Riot}, Regulation, Winter 1995, at 83, 84；Barnett）。

\subsection{（三）審慎根據是自制僅餘的理由\pg{633 頁}}

Kramer\ix{persons}{kramer@Kramer, Larry D.} 說，在兩派，推動新解釋理論\ix{terms}{05@\zhlidxcat{解釋理論與自制}!解釋理論}發展的同一個變化，即接受終局的解釋權在司法，也削弱了持久的尊重\ix{terms}{05@\zhlidxcat{解釋理論與自制}!尊重}或自制\ix{terms}{05@\zhlidxcat{解釋理論與自制}!司法自制}學說的基礎。穩固的尊重學說的基礎，在於有人可以尊重，且有好理由。接受司法至上\ix{terms}{01@\zhlidxcat{司法審查權限之歸屬}!司法至上}，消除了原本 Jeffersonian\ix{terms}{01@\zhlidxcat{司法審查權限之歸屬}!共和派} 的尊重基礎（尊重並保存美國人民首要的解釋權威），也消除了其 Thayer\ix{persons}{thayer@Thayer, James Bradley} 式的後繼者（把該權威轉給被認為最能回應人民的政治部門）。

剩下的只有審慎\ix{terms}{05@\zhlidxcat{解釋理論與自制}!審慎}的根據：不要因行動過於頻繁或積極而引發政治反撲，以確保司法權威。Kramer 的問句是：「Why should the Court defer, after all, if it is, by constitutional design, supposed to be the final, authoritative interpreter?」〔法院若依憲法設計本應是終局的、權威的解釋者，究竟為何應該尊重？〕答案是除了保護法院扮演其制度角色的能力之外沒有理由。在司法至上的世界，尊重「simply a matter of expediency」〔僅是權宜問題〕，即法院在實務上可走多遠而不致耗掉太多政治資本。據此，左右兩派都要求法院尊重，但只限於政治過程的結果是他們贊成或不太在意的領域，Kramer 稱「hardly surprising」〔不足為奇〕\pg{633 頁}。

\subsection{（四）均衡點\pg{633 至 634 頁}}

司法至上\ix{terms}{01@\zhlidxcat{司法審查權限之歸屬}!司法至上}愈深入國家的政治文化，感到需要尊重\ix{terms}{05@\zhlidxcat{解釋理論與自制}!尊重}的程度就愈低\pg{633 至 634 頁}。Kramer\ix{persons}{kramer@Kramer, Larry D.} 說，任何容許司法審查\ix{terms}{01@\zhlidxcat{司法審查權限之歸屬}!司法審查}的制度，即使採司法至上，都有一個均衡點\ix{terms}{05@\zhlidxcat{解釋理論與自制}!均衡點}，法院超過它就會損害其制度權威與行動能力。人們對司法審查的反應也受人們對法院應有多少權威之感受的影響，不單取決於是否同意法院的結果：「People will defer to a superior with whom they disagree more than a subordinate」〔人們對上位者的不同意見\ix{terms}{04@\zhlidxcat{法院之功能}!不同意見與協同意見}，比對下屬更能容忍〕，他舉的例子是院長做的事教師會接受，職員做就不會\pg{634 頁}。因此，若人們相信終局決定權屬於法院，自己與代表都不握有該權，就會讓大法官走得更遠、做更多自己不喜歡的事。

Kramer 稱「Judicial supremacy is an ideology」〔司法至上是一種意識形態〕，其全部目的與效果，是把公眾與政治接受的均衡點移向有利於司法權威的一側。當司法至上在整個政治體系中被視為理所當然，法院不必顧慮權威即可行動的能力就自然增長，尊重所起的作用也愈來愈小，直到如 Posner\ix{persons}{posner@Posner, Richard A.} 所言，成為「little more than a rhetorical tool used opportunistically by pretty much all of the Justices」〔幾乎所有大法官都機會主義地使用的修辭工具而已〕\pg{634 頁}。

\subsection{（五）結論\pg{634 頁}}

全文最後一句，Kramer\ix{persons}{kramer@Kramer, Larry D.} 說若這是壞事，而他認為是，禍首在「our embrace of the strange and counterintuitive notion that determining these answers, and so deciding the ongoing meaning of a democratic Constitution, is best left to life-tenured judges whose chief qualification is technical legal training」〔我們接受了這個奇怪而違反直覺的想法：決定這些答案，因而決定民主憲法持續演變的意義，最宜交給終身任職、主要資格是技術性法律訓練的法官〕。同一句開頭已把解釋理論\ix{terms}{05@\zhlidxcat{解釋理論與自制}!解釋理論}，以及憲法問題有正確答案這個想法，排除在禍首之外。
